\chapter{Dataset and Experimental Design}

\section{Synthetic Dataset Generation}

The dataset is generated by randomly sampling option and market parameters from fixed ranges.
Each observation initially contains the primitive Black-Scholes input variables:

\begin{equation}
    x = (S_0, K, T, r, \sigma),
\end{equation}

and the target is:

\begin{equation}
    y = C_{BS}(S_0, K, T, r, \sigma).
\end{equation}

Using synthetic data is a methodological choice.
The project studies whether a neural network can approximate a known mathematical pricing function, not whether it can fit noisy market option prices.
With synthetic Black-Scholes data, the ground truth is exact and reproducible.

\section{Input Variables and Target}

The parameter ranges used to generate the dataset are reported in Table~\ref{tab:dataset-ranges}.

\begin{table}[H]
    \centering
    \begin{tabular}{lll}
        \toprule
        Variable & Description              & Range        \\
        \midrule
        $S_0$    & initial underlying price & $[50, 150]$  \\
        $K$      & strike price             & $[50, 150]$  \\
        $T$      & maturity in years        & $[0.1, 2]$   \\
        $r$      & risk-free rate           & $[0, 0.05]$  \\
        $\sigma$ & volatility               & $[0.1, 0.6]$ \\
        \bottomrule
    \end{tabular}
    \caption{Parameter ranges used to generate the synthetic dataset.}
    \label{tab:dataset-ranges}
\end{table}

The target variable is the European call option price computed using the analytical Black-Scholes formula.
Therefore, each supervised learning observation has the form:

\begin{equation}
    (S_0, K, T, r, \sigma) \mapsto C_{BS}.
\end{equation}

\section{Feature Engineering: Moneyness}

The final model augments the input vector with moneyness:

\begin{equation}
    \mathrm{moneyness} = \frac{S_0}{K}.
\end{equation}

The final input vector is:

\begin{equation}
    x_{\mathrm{final}} = (S_0, K, T, r, \sigma, S_0/K).
\end{equation}

Moneyness does not add new information beyond $S_0$ and $K$, but it exposes a financially meaningful relationship directly to the model.
Intermediate experiments showed that this engineered feature can improve absolute-error metrics.

\section{Train/Validation/Test Split}

The dataset is split into training, validation, and test sets.
The training set is used to optimize the model parameters.
The validation set is used for early stopping and model selection during training.
The test set is kept separate and is used only for final evaluation.

All splits are controlled by fixed random seeds.
This makes the reported results reproducible and ensures that different model variants can be compared under consistent conditions.

\section{Clean and Noisy Target Settings}

The main experiment uses clean Black-Scholes targets.
The project also includes a controlled noisy-target setting used for robustness analysis.
In that setting, the input variables are unchanged, while the training labels are perturbed.

For each clean price $C_{\mathrm{BS}}$, the noisy target is generated as:

\begin{equation}
    \widetilde{C}
    =
    \max\left(
    C_{\mathrm{BS}} + \varepsilon,
    0
    \right),
    \qquad
    \varepsilon \sim
    \mathcal{N}\left(
    0,
    \left(\ell \max(C_{\mathrm{BS}}, 1)\right)^2
    \right),
\end{equation}

where $\ell$ is the noise level.
The evaluation target remains the clean analytical Black-Scholes price.
This experiment measures robustness to imperfect labels while preserving a known ground truth.

\section{Reproducibility}

Reproducibility is handled through fixed random seeds and explicit configuration objects.
Command-line options define each run, and the corresponding artifacts are saved to disk.
Generated datasets and model checkpoints are stored in reproducible output directories.
Selected final metrics and figures are tracked in the repository.
